%============================== MÉTHODOLOGIE ================================
\chapter{Méthodologie de travail}

\section{Une démarche itérative et hebdomadaire}

Le développement a suivi une \textbf{démarche itérative}, rythmée par des
\textbf{cycles hebdomadaires}. Chaque semaine ciblait un ou plusieurs modules
cohérents, développés puis validés avant de passer aux suivants. Cette
approche présente plusieurs avantages : elle limite les risques (un module
défaillant est isolé), elle rend l'avancement \textbf{traçable}, et elle permet
un retour régulier avec l'encadrante.

Le tableau~\ref{tab:planning} résume la progression des modules livrés au fil
des semaines.

\begin{table}[H]
\centering
\caption{Progression itérative du projet}
\label{tab:planning}
\begin{tabularx}{\textwidth}{@{}p{2.6cm}L@{}}
\toprule
\textbf{Étape} & \textbf{Modules / travaux principaux} \\
\midrule
Semaine 2 & Fondation frontend : structure, authentification simulée, rôles, layout, tableau de bord initial. \\
Semaine 3 & Assistant de création de campagne (formulaire multi-étapes). \\
Semaine 4 & Import de la base de magasins, carte interactive, sélection et \emph{geofencing}, intégration au parcours de campagne. \\
Semaine 5 & DCO : upload de visuels, génération de variantes, \emph{landing pages} simulées ; reporting (indicateurs, graphiques, carte, tableau, export CSV). \\
Semaine 6 & Gestion du compte (annonceurs, utilisateurs, paramètres), intégration de bout en bout, corrections, documentation technique, préparation et livraison de la démonstration finale. \\
\bottomrule
\end{tabularx}
\end{table}

\section{Flux de travail Git / GitHub}

Le code et la documentation sont versionnés avec \textbf{Git} et hébergés sur
\textbf{GitHub}. Le flux de travail repose sur :

\begin{itemize}
  \item une \textbf{branche par sujet} (fonctionnalité, correctif ou
        documentation), nommée de façon explicite (par exemple
        \texttt{feat/week6-account-management},
        \texttt{docs/week6-technical-documentation},
        \texttt{docs/final-report}) ;
  \item des \textbf{\emph{pull requests}} pour intégrer chaque branche, ce qui
        matérialise chaque incrément et conserve un historique lisible ;
  \item des \textbf{messages de commit} et des libellés de PR décrivant le
        contenu de chaque étape.
\end{itemize}

La figure~\ref{fig:github} présente le dépôt GitHub du projet et son
historique.

\section{Validation après chaque module}

Chaque module a fait l'objet d'une \textbf{validation fonctionnelle} avant
d'être considéré comme terminé :

\begin{itemize}
  \item vérification du \textbf{build} du frontend (\texttt{npm run build}) et
        de l'analyse statique (\texttt{npm run lint}) ;
  \item vérification de la \textbf{santé du backend}
        (\texttt{GET /api/health}) et des endpoints via \textbf{Swagger} ;
  \item \textbf{parcours manuel} dans le navigateur (navigation, formulaires,
        cas d'erreur) ;
  \item contrôle de \textbf{non-régression} sur les modules déjà validés.
\end{itemize}

Le détail de ces validations est consigné dans le journal d'avancement du
dépôt (fichiers \texttt{docs/etat-avancement-semaine-*.md}).

\section{Documentation continue et préparation de la démonstration}

La documentation a été produite \textbf{en continu}, et non reléguée en fin de
projet. Le dépôt contient notamment une documentation technique complète, un
schéma de base de données, un scénario de démonstration, une checklist de
démonstration et un document de livraison finale. Cette production
documentaire fait l'objet d'un chapitre dédié.
